\section{The period curvature of the retained critical-comparison kernel}
\label{sec:period-critical-kernel-bridge}

This section joins the original theta representatives, the critical-line
comparison of the Hochschild note (H16)--(H25), and the period curvature
(SC1)--(SC37).  The single-packet identification is proved first, including
its full arithmetic unit.  We then calculate the factorwise critical
comparison on every cyclic tensor packet.  All tensor targets below use
the completed projective tensor product of the programme's locally convex
spaces; its finite surviving image is proved injective by a continuous
retraction, without an assertion of tensor exactness for arbitrary spaces.

The proof-bearing comparison source is
\emph{The original theta complex inside the Hochschild trace realization},
(H1)--(H25), at the retained PR27 revision
\nolinkurl{a4a87844c257485ae6d669f5b3eb6d96b9e8b8d8}; its original
analytic constructions are retained in the cumulative source.
The arithmetic section and unit maps are (A7)--(A13), (CT.1)--(CT.5),
and (TVB.1)--(TVB.3).  The calculations here prove the additional maps
between those specified objects.

\subsection{The single-packet polynomial and its complete arithmetic unit}

Retain the original nonempty packet, with every selected zero's complete
order,
\[
 h(s)=\prod_{\rho\in Z}(s-\rho)^{m_\rho},\qquad
 g(s)=2\xi(s),\quad v_h(s)=g(s)/h(s),\quad d=\deg h.
 \tag{PC1}
\]
The packet is the same one used by the original source, including its
reflection symmetry when that symmetry is used in its metric.  Put
$E_h=\mathbb C[s]/(h)$, and let $j_h$ mean the full Hermite remainder.
Its local coordinate at $\rho$ is $x=s-\rho$, with basis
$1,x,\ldots,x^{m_\rho-1}$.  Thus
\[
 \upsilon_{h,\rho}
   =\sum_{a=0}^{m_\rho-1}\frac{v_h^{(a)}(\rho)}{a!}x^a,
 \quad\upsilon_h=j_h(v_h),\quad
 \varepsilon_h=\upsilon_h^{-1}.
 \tag{PC2}
\]
This inverse exists with all its coefficients retained.  Indeed write
$h(s)=(s-\rho)^{m_\rho}h_\rho(s)$ and
$g(s)=(s-\rho)^{m_\rho}g_\rho(s)$.  Both remaining factors are nonzero
at $\rho$, so the constant coefficient
$a_0=g_\rho(\rho)/h_\rho(\rho)$ is nonzero.  If the coefficients
of $\upsilon_{h,\rho}$ are $a_a$, its inverse coefficients satisfy
\[
 b_0=a_0^{-1},\qquad
 b_n=-a_0^{-1}\sum_{a=1}^n a_a b_{n-a}\quad(1\le n<m_\rho).
 \tag{PC3}
\]
Multiplication makes the constant coefficient one and each other
coefficient through degree $m_\rho-1$ zero.  The coprime polynomial
factors of $h$ give the Chinese-remainder isomorphism; assembling these
inverses proves (PC2).  The notation $j_h(h/g)$ for $\varepsilon_h$
uses its holomorphic germ at every selected zero; no global entire
reciprocal is asserted.

For $k=1$, the sum variable is literally $S=s_1$.  The isomorphism
$\mathbb C[S]\to\mathbb C[s_1]$, $S\mapsto s_1$, pulls
$(h(s_1))$ back to $(h(S))$.  The monic annihilator defining the
cyclic packet is consequently
\[
 \begin{gathered}
 \chi_{h,1}(S)=h(S),\quad E=\mathbb C[S]/(h(S)),\quad
 A=M_S,\quad q=d,\\
 \iota_1:E\xrightarrow{\sim}E_h,\quad[P(S)]\mapsto[P(s)],
 \quad\eta_h=M_{\upsilon_h}\iota_1.
 \end{gathered}
 \tag{PC4}
\]
The equality includes every exponent $m_\rho$.  The inverse of
$\eta_h$ is $\iota_1^{-1}M_{\varepsilon_h}$.  It intertwines the
two original multiplication actions because multiplication by the unit
commutes with $s$.  Its actual multiplicative rule is
$\eta_h(x)\eta_h(y)=\upsilon_h\eta_h(xy)$ and
$\eta_h(1)=\upsilon_h$.  Thus its source and target polynomial
algebras and its module-map type remain explicit.

Use the original spaces and maps
\[
\begin{gathered}
 V=\{\phi\in\mathcal S(\mathbb R)_{\rm ev}:
             \phi(0)=\int_{\mathbb R}\phi=0\},\quad
 D=-x\partial_x,\quad\Theta\phi(x)=2\sum_{n\ge1}\phi(nx),\\
 \mathscr B=\{F\in C^\infty(\mathbb R_{>0}):\sup_{x>0}x^b|D^jF(x)|<\infty
                  \text{ for all }b\in\mathbb Z,j\ge0\},\quad
 Q=\mathscr B/\Theta V,\quad\mathsf q:\mathscr B\to Q,\\
 \mathcal MF(s)=\int_0^\infty F(x)x^{s-1}\,dx,\quad
 \phi_*=(4\pi^2x^4-6\pi x^2)e^{-\pi x^2},\quad
 \mathcal M F_h=v_h,\quad h(D)F_h=\Theta\phi_*.
\end{gathered}\tag{PC5}
\]
The source construction of $F_h$ is the iterated original Euler inverse
$S_\rho F=x^{-\rho}\int_x^\infty F(y)y^{\rho-1}dy$ on the
vanishing Mellin fibre, as proved in (A5)--(A7) and (TVB.1).
Let $\mathcal T_hP=P(D)F_h$, $\pi_hP=[P]_h$, and
$J_h=j_h\mathcal M:\mathscr B\to E_h$, with the retained
$\sigma_h:E_h\to Q$.  The exact source relations are
\[
 \mathcal T_h(hP)=\Theta(P(D)\phi_*),\qquad
 J_h\mathcal T_h=\eta_h\pi_h,\qquad
 \mathsf q\mathcal T_h=\sigma_h\eta_h\pi_h.
 \tag{PC6}
\]
Here and below $\pi_h$ uses the variable identification (PC4).
The first identity follows from $h(D)F_h=\Theta\phi_*$ and
commutation with $D$.  Integration by parts gives
$\mathcal M\mathcal T_hP=v_hP$, proving the second identity.
For the third, the retained section in (A10) is
$\sigma_h=\mathsf q R_{\rm ref}$, with
$R_{\rm ref}u=\mathcal T_h\operatorname{rem}_h(\varepsilon_hu)$.
Write $P=P_0+hT$, $\deg P_0<d$.  Then
$R_{\rm ref}\eta_h[P]=\mathcal T_hP_0$, and the difference from
$\mathcal T_hP$ is the displayed theta boundary
$\Theta(T(D)\phi_*)$.  This proves the precise equality with that
section, rather than identifying two sections merely from equal jets.
Since $J_h\Theta=0$, its descended map is
$\overline J_h:Q\to E_h$, defined by
$\overline J_h\mathsf q=J_h$; the section identity is
$\overline J_h\sigma_h=I_{E_h}$.
Also $D\mathcal T_hP=\mathcal T_h(SP)$ and (PC6) imply
$D_Q\sigma_h\eta_h=\sigma_h\eta_hA$.
Surjectivity of $\eta_h$ therefore proves
$D_Q\sigma_h=\sigma_hM_s$ with its exact domains.

For each admitted original polynomial degree $N$, retain its actual
remainder matrix $B_N$, source Gram $O_N$, and least representative
\[
 K_N=B_NO_N^{-1}B_N^*,\quad G_N=K_N^{-1},\quad
 R_N=O_N^{-1}B_N^*G_N,\quad r_N=\mathcal T_hR_N.
 \tag{PC7}
\]
The source Gram is for the original density
$|v_h(1/2+iy)|^2dy/(2\pi)$ and the original monic coordinates.
Since $B_NR_N=I$, (PC6) gives
\[
 j_E:=\eta_h^{-1}J_h:\mathscr B\to E,\qquad
 j_Er_N=I,\quad J_hr_N=\eta_h,\quad
 \mathsf q r_N=\sigma_h\eta_h,\quad j_ED=Aj_E.
 \tag{PC8}
\]
The last identity uses $J_hD=M_sJ_h$ and the intertwining in
(PC4).  In particular the coefficient observer in (SC7) is $j_E$.
In full Mellin-jet coordinates the same Gram is
$\eta_h^{-*}G_N\eta_h^{-1}$; multiplying back by $\eta_h^*$
and $\eta_h$ recovers exactly $G_N$.  The unit is present in both
the source map and the coordinate congruence.

\subsection{The weighted critical comparison and every recovered jet}

Reserve $y\in\mathbb R$ for the critical-line coordinate, distinct
from the complex period parameter $t$.  The retained continuous map is
\[
 Q_{\rm crit}=\mathcal S(\mathbb R)\big/
   \overline{\zeta(1/2+iy)\mathcal S(\mathbb R)},\qquad
 \Gamma[F]=[\mathcal MF(1/2+iy)],\qquad
 M=M_{1/2+iy},\quad \Gamma D_Q=M\Gamma.
 \tag{PC9}
\]
It is induced by the plus-sign Mellin restriction: the theta image
restricts to $2\zeta(1/2+iy)\mathcal M\phi(1/2+iy)$.
The continuity and strong-Schwartz domains are those proved in
(H16)--(H20); continuity also makes it factor through the closed
theta quotient.  Let $Z_c=\{\rho\in Z:\Re\rho=1/2\}$ and
$Z_o=Z\setminus Z_c$.  Retain the complete factors
\[
 h_c=\prod_{\rho\in Z_c}(S-\rho)^{m_\rho},\quad
 h_o=\prod_{\rho\in Z_o}(S-\rho)^{m_\rho},\quad h=h_ch_o.
 \tag{PC10}
\]
An empty product is $1$.  Chinese remainders identify
$E=E_c\oplus E_o$, where $E_c$ is the ideal $(h_o)/(h)$ and
$E_o$ is the ideal $(h_c)/(h)$, with projections $\pi_c,\pi_o$.
The notation also retains the isomorphisms to the products of their
local primary algebras, including their nilpotent coordinates.

Define the actual weighted comparison
\[
 \widetilde\sigma_h=\sigma_h\eta_h:E\to Q,\qquad
 \mathcal C_h=\Gamma\widetilde\sigma_h:E\to Q_{\rm crit}.
 \tag{PC11}
\]
Then
\[
 \boxed{\ker\mathcal C_h=E_o,\qquad
         \mathcal C_hA=M\mathcal C_h.}
 \tag{PC12}
\]
To prove the kernel, first take an off-critical root $\rho$.
The reciprocal $(1/2+iy-\rho)^{-1}$ and all its derivatives have
polynomial bounds on the real axis.  Multiplication by it preserves
Schwartz space, the zeta ideal and its closure, and is inverse to
$M-\rho$ on $Q_{\rm crit}$.  On a primary block
$(A-\rho)^{m_\rho}=0$, intertwining therefore forces the whole
block to zero.  Multiplication by $\upsilon_h$ preserves each CRT
block and does not change this conclusion.

For a critical root $\rho=1/2+i\gamma$ retain the evaluation
\[
 \Xi_\rho[f]=\sum_{j=0}^{m_\rho-1}
       \frac{i^{-j}}{j!}f^{(j)}(\gamma)x^j
       \in\mathbb C[x]/(x^{m_\rho}),\qquad x=S-\rho.
 \tag{PC13}
\]
Every derivative used here vanishes on the zeta ideal and, by
continuity, on its closure.  The chain rule supplies $i^j$ when
differentiating $\mathcal MF(1/2+iy)$, so (PC6) and
$\overline J_h\sigma_h=I$ give the full formula
\[
 \Xi_\rho\mathcal C_h[P]
 =\upsilon_{h,\rho}[P]_\rho
 =\sum_{j<m_\rho}\sum_{a=0}^j
   \frac{v_h^{(a)}(\rho)}{a!}
   \frac{P^{(j-a)}(\rho)}{(j-a)!}x^j.
 \tag{PC14}
\]
In particular the critical observer recovers the actual unit on
$e_0=[1]$, rather than the constant one.  Assemble the $\Xi_\rho$
by the retained CRT map and then $\iota_1^{-1}$ into
$\Xi_c:Q_{\rm crit}\to E_c$.  Put
$\upsilon_c=\pi_c\iota_1^{-1}(\upsilon_h)$ and
$\varepsilon_c=\pi_c\iota_1^{-1}(\varepsilon_h)$; these are
inverse units in the ideal algebra $E_c$, whose identity is
$e_c=\pi_c(e_0)$.  Multiplication by this inverse defines the continuous map
\[
 \mathcal R_c:=M_{\varepsilon_c}\Xi_c,
 \qquad \mathcal R_c\mathcal C_h=\pi_c.
 \tag{PC15}
\]
This is a left inverse on the entire critical block, proving (PC12).
For an empty critical packet both maps have zero-dimensional target.
When $E_c\ne0$, (PC14) proves $\mathcal C_h(e_0)\ne0$.
The finite image $W_h=\mathcal C_h(E_c)$ is preserved by $M$,
and $\mathcal C_h|E_c$ and $\mathcal R_c|W_h$ are inverse maps.

\subsection{The comparison kernel is a holomorphic period subbundle}

Use exactly the polynomial $\chi=h$ in (SC1)--(SC6), its potential
\[
 \Phi_t(S)=\frac{S^{q+1}}{q+1}
   +\sum_{a=0}^{q-1}\frac{c_aS^{a+1}}{a+1}-tS,
 \quad h(S)=S^q+\sum_{a=0}^{q-1}c_aS^a,
 \quad u\in\mathbb C^*.
 \tag{PC16}
\]
Fix the original argument of $u$ and rays
$\theta_j=(\arg u+\pi+2\pi j)/(q+1)$, with
$\Gamma_j=-\ell_0^{\rm ray}+\ell_j^{\rm ray}$ and
$\Pi_{jb}(t)=\int_{\Gamma_j}e^{\Phi_t(S)/u}S^b\,dS$.
The original monomial factors are
\[
 (\Pi_{\rm mon})_{j,b}=
 \frac{((q+1)|u|)^{(b+1)/(q+1)}}{q+1}
 \Gamma\!\left(\frac{b+1}{q+1}\right)e^{i(b+1)\theta_0}
 (e^{2\pi i j(b+1)/(q+1)}-1).
 \tag{PC17}
\]
They have nonzero determinant: a column dependence of their last
matrix gives a polynomial of degree at most $q$ vanishing at all
$q+1$ roots of unity.  Its coefficients must be zero.  Along any
compact coefficient path the negative leading ray term dominates
the lower-degree terms by an integrable bound
$Cr^L\exp(-r^{q+1}/(2(q+1)|u|))$.  Differentiation under the
integral is therefore valid.  The original polynomial reduction
$uP'+(h-t)P$ has leading degree $\deg P+q$ and the same leading
coefficient, so it has a unique remainder below degree $q$.
Its integral vanishes by the exact derivative and cancellation of
the two origin values.  Coefficient variation thus gives
$\Pi'=\Pi B$ along a path, and the adjugate determinant identity
gives $\det\Pi(s)=\det\Pi(0)\exp(\int\operatorname{Tr}B)$.
This proves invertibility at every $t$, including all collisions.

In the unchanged coefficient basis set
\[
 \ell([P])=[S^{q-1}]\operatorname{rem}_hP,\quad
 \mathcal R=e_0\ell,\quad A_t=A+t\mathcal R,\quad
 \Pi'(t)=-\Pi(t)A_t/u.
 \tag{PC18}
\]
The last equality follows by differentiating in $t$: for the final
column the sole reduction is $[h-t]=0$, obtained from the
differential of the constant polynomial one.
Define maps on the period target, with range still in the original
critical quotient,
\[
 \mathcal C_{h,t}=\mathcal C_h\Pi(t)^{-1},\qquad
 A_t^{\rm per}=\Pi(t)A_t\Pi(t)^{-1},\qquad
 \mathcal K_{h,t}=\ker\mathcal C_{h,t}=\Pi(t)E_o.
 \tag{PC19}
\]
The equality of kernels follows from invertibility and (PC12).
The inverse matrix is holomorphic, so this is a holomorphic kernel
bundle of constant rank, with the indicated global frame from $E_o$.
Its quotient is the fixed finite image $W_h$; more explicitly
$\mathcal C_{h,t}\Pi(t)|E_c=\mathcal C_h|E_c$ and
$\Pi(t)|E_c\,\mathcal R_c$ is a right inverse on $W_h$.

The exact action defect is
\[
 \boxed{\mathcal C_{h,t}A_t^{\rm per}-M\mathcal C_{h,t}
       =t\,\mathcal C_h(e_0)\ell\Pi(t)^{-1}.}
 \tag{PC20}
\]
Indeed multiplying the first term gives
$\mathcal C_h(A+t\mathcal R)\Pi^{-1}$, and the part containing
$A$ cancels $M\mathcal C_h\Pi^{-1}$ by (PC12).  This proves
the formula before restricting its domain.

Put $F=E_o$, $p=\dim F$.  If $0<p<q$, it is the ideal
$(h_c)/(h)$, with ordered basis
$[h_c],[Sh_c],\ldots,[S^{p-1}h_c]$.
The last member is monic of degree $q-1$, so
$\ell|F$ is onto with kernel dimension $p-1$.
An ideal containing $e_0$ would be all of $E$, hence $e_0\notin F$.
It follows from (PC14) and (PC20) that, for $t\ne0$, the
restriction of this action defect to $\mathcal K_{h,t}$ has
\[
 \operatorname{im}=\mathbb C\mathcal C_h(e_0),\quad
 \ker=\Pi(t)(F\cap\ker\ell),\quad\operatorname{rank}=1.
 \tag{PC21}
\]
At $t=0$ it is zero.  If $F=0$ its restricted domain is zero.
If $F=E$, the comparison $\mathcal C_h$ is zero and every
displayed defect is zero.  These statements cover all actual
sector possibilities without asserting that an off-critical zero exists.

\subsection{The full curvature and both original metric observations}

Take any specified coefficient inclusion $I:F\hookrightarrow E$;
put $L=\ell I$, and for $0<p<q$ define
\[
 Y=\Pi I,\quad H=Y^*Y,\quad
 P=YH^{-1}Y^*,\quad N=I_q-P,\quad z=N\Pi e_0.
 \tag{PC22}
\]
This is exactly the kernel bundle (PC19), measured in the fixed
period-coordinate Gram $I_q$.  The Gram is positive because $Y$
is injective.  Multiplication proves that $P$ is its orthogonal
projection.  Since $e_0\notin F$ and $\Pi$ is invertible,
$z(t)\ne0$ for every $t$.  Original invariance $AI=IA_F$ gives
\[
 Y'=-YA_F/u-(t/u)\Pi e_0L,
 \quad NY'=-(t/u)zL,
 \quad N(0)Y''(0)=-z(0)L/u.
 \tag{PC23}
\]
The final equation follows by differentiating (PC18); the term
$\Pi(0)A^2I/u^2$ lies inside the kernel frame and is removed by
$N(0)$.  Thus the first normal derivative vanishes at zero and
the normal acceleration has rank one, with its original $u$ factor.

For $t=x+iy_t$ use
$\partial_t=(\partial_x-i\partial_{y_t})/2$ and its conjugate.
Holomorphicity gives $\partial_tH=Y^*Y'$,
$\partial_{\bar t}H=Y'^*Y$, and
$\partial_{\bar t}\partial_tH=Y'^*Y'$.
Differentiate $H^{-1}H=I$ and the determinant to obtain
\[
\begin{split}
 \partial_{\bar t}\partial_t\log\det H
 &=\operatorname{Tr}(H^{-1}Y'^*Y'
     -H^{-1}Y'^*YH^{-1}Y^*Y')\\
 &=\operatorname{Tr}(H^{-1}Y'^*NY')
   =\frac{|t|^2}{|u|^2}\|z\|^2 L H^{-1}L^*.
\end{split}\tag{PC24}
\]
Every factor in the last coefficient is strictly positive except the
explicit $|t|^2$: $z\ne0$, $L\ne0$ and $H^{-1}>0$.
The coefficient is real analytic, so the curvature vanishes at zero
to exactly quadratic order.  In the determinant-dual metric
$(\det H)^{-1}$ its first Chern form is
$\frac{i}{2\pi}\frac{|t|^2}{|u|^2}\|z\|^2LH^{-1}L^*
dt\wedge d\bar t$, with $i\,dt\wedge d\bar t=2dx\wedge dy_t$.
This proves that the SC curvature is the curvature of this actual
critical-comparison kernel in the specified fixed period metric.

The original theta Gram remains separately and explicitly present.
Let $G_{N,F}=I^*G_NI$, and set
\[
 B_{N,F}=G_{N,F}^{-1}H,\quad
 \widetilde G_N=\Pi^{-*}G_N\Pi^{-1},\quad
 P_N^\theta=\Pi I G_{N,F}^{-1}I^*G_N\Pi^{-1}.
 \tag{PC25}
\]
Thus $H=G_{N,F}B_{N,F}$ and
$Y^*\widetilde G_NY=G_{N,F}$.  Both identities follow by direct
multiplication.  The squared normal norm of $\Pi e_0$ for the
moving target metric is exactly
\[
 e_0^*G_Ne_0-e_0^*G_NI G_{N,F}^{-1}I^*G_Ne_0>0.
 \tag{PC26}
\]
It is positive because it is the squared norm of the nonzero
$G_N$-orthogonal component of $e_0$ outside $F$.  From (PC18),
$\partial_t\Pi^{-1}=A_t\Pi^{-1}/u$, so
\[
 \partial_t\widetilde G_N=\Pi^{-*}G_NA_t\Pi^{-1}/u,
 \qquad
 \partial_{\bar t}\widetilde G_N=
       \Pi^{-*}A_t^*G_N\Pi^{-1}/\bar u.
 \tag{PC27}
\]
In the product rule for $Y^*\widetilde G_NY$ these terms cancel
the derivatives of $Y$ and $Y^*$ exactly.  Hence its Gram and
determinant are constant in the same holomorphic frame.  Formula
(PC24) and this cancellation are related by the full displayed
metric congruence, including all original theta masses and unit factors.
For $F=0$ the determinant is $1$; for $F=E$, $N=0$.
Both endpoint curvatures are therefore zero.

\subsection{The exact completed tensor comparison and its collision jets}

We now retain an arbitrary original integer $k\ge1$.  Set
\[
 B_k=E_h^{\otimes k}
  =\mathbb C[s_1,\ldots,s_k]/(h(s_1),\ldots,h(s_k)),\quad
 S=s_1+\cdots+s_k,\quad U_k=\upsilon_h^{\otimes k},
 \tag{PC28}
\]
and retain the original cyclic algebra and its two distinct maps.
The variable isomorphism between its tensor coefficient carriers is
$\iota_k=\iota_1^{\otimes k}:E^{\otimes k}\xrightarrow{\sim}B_k$.
Thus
\[
 E_k=\mathbb C[S]/(\chi_{h,k}),\quad A_k=M_S,\quad
 \alpha_k[P]=[P(S)]\in B_k^{\mathfrak S_k},\quad
 \eta_k=M_{U_k}\alpha_k.
 \tag{PC29}
\]
The map $\alpha_k$ is unital; $\eta_k$ is the original unit-weighted
module injection.  Define the fully typed comparison
$\overline\alpha_k=\iota_k^{-1}\alpha_k:E_k\to E^{\otimes k}$.
Then $\eta_k=\eta_h^{\otimes k}\overline\alpha_k$, since
$\eta_h=M_{\upsilon_h}\iota_1$.  Their source polynomial is exactly
\[
 \chi_{h,k}(S)=\prod_\lambda(S-\lambda)^{m_\lambda},\quad
 m_\lambda=\max_{\substack{\rho_1+\cdots+\rho_k=\lambda\\\rho_i\in Z}}
          \left(1+\sum_{i=1}^k(m_{\rho_i}-1)\right).
 \tag{PC30}
\]
All tuples are ordered in the full tensor algebra; equal sums are
grouped by their literal equality.  To verify the exponent, the local
factor at a tuple has coordinates $x_i=s_i-\rho_i$ and quotient
by $(x_i^{m_{\rho_i}})$.  The nilpotent sum
$x_1+\cdots+x_k$ has highest nonzero power
$D=\sum_i(m_{\rho_i}-1)$: that power contains the monomial
$\prod_i x_i^{m_{\rho_i}-1}$ with the nonzero coefficient
$D!/\prod_i(m_{\rho_i}-1)!$.  Higher powers vanish.  Taking the
maximum over equal sums proves (PC30) and the injectivity of
$\alpha_k$.  The invariant unit vector already has that same
annihilator, so passage to symmetric invariants does not alter it.

Let $\chi_{c,k}$ be the same polynomial formed using only tuples
from $Z_c$, with exponents $r_\lambda$ given by the same maximum;
an absent critical tuple at $\lambda$ has $r_\lambda=0$.
When $Z_c=\varnothing$, define $\chi_{c,k}=1$.  Put
$q_k=\deg\chi_{h,k}$ and $r_k=\deg\chi_{c,k}$.
The tuple inclusion proves $r_\lambda\le m_\lambda$ and hence
$\chi_{c,k}\mid\chi_{h,k}$.  The same nilpotence proof gives
the exact injection
\[
 \alpha_{c,k}:\mathbb C[S]/(\chi_{c,k})
                  \hookrightarrow E_c^{\otimes k},\quad
 [P]\mapsto P(A_c^{[k]})e_c^{\otimes k},
 \quad
 \pi_c^{\otimes k}\overline\alpha_k=
          \alpha_{c,k}\pi_{\chi_{c,k}}.
 \tag{PC31}
\]
Here $A_c=A|E_c$ and $A_c^{[k]}$ is its sum on the ordered
tensor factors.  The identity vector is the displayed $e_c^{\otimes k}$,
so the codomain is precisely the original CRT ideal tensor.
For an empty critical factor both the source and target of
$\alpha_{c,k}$ are zero.  This identity keeps the full source
polynomial $\chi_{h,k}$ and the actual quotient map to its critical
tuple polynomial.

The target is the completed projective tensor space
$Q_{\rm crit}^{\widehat\otimes_\pi k}$.  There is no loss of the
finite critical image on completion.  In fact define
$i_c=\mathcal C_h|E_c$ and use the continuous map $\mathcal R_c$
in (PC15).  Their tensor maps extend continuously to the projective
completions, and on pure tensors their composition is identity.
Pure tensors span the finite-dimensional $E_c^{\otimes k}$, which
is already complete, so
\[
 \mathcal R_c^{\widehat\otimes k}
             i_c^{\widehat\otimes k}=I_{E_c^{\otimes k}}.
 \tag{PC32}
\]
This proves injectivity and a continuous retraction onto the finite
image inside the completed target.  Its associated continuous
idempotent has that image as its kernel-complement, so the image is
closed in the Hausdorff completed space.  This argument uses the
displayed maps and does not assume preservation of arbitrary exact
sequences by completion.

The actual tensor comparison, with its original weighted inclusion, is
\[
 \mathcal C_k=
    \Gamma^{\widehat\otimes k}\sigma_h^{\widehat\otimes k}\eta_k
    =\mathcal C_h^{\widehat\otimes k}\overline\alpha_k:
        E_k\to Q_{\rm crit}^{\widehat\otimes_\pi k}.
 \tag{PC33}
\]
Here $\sigma_h^{\widehat\otimes k}\eta_k$ is the existing arithmetic
cohomology inclusion.  Applying the raw critical evaluations gives
$\Xi_c^{\widehat\otimes k}\mathcal C_k=
(M_{\upsilon_c})^{\otimes k}\alpha_{c,k}\pi_{\chi_{c,k}}$.
Applying their inverses (PC32) instead gives
$\alpha_{c,k}\pi_{\chi_{c,k}}$.  Every derivative coordinate has
its product phase $i^{-\sum j_i}/\prod j_i!$ and the full tensor
of the local units (PC14).  The injectivity in (PC31)--(PC32)
therefore proves the exact kernel and image dimension
\[
 \boxed{F_k:=\ker\mathcal C_k=(\chi_{c,k})/(\chi_{h,k}),
       \qquad\dim F_k=q_k-r_k,\quad
       \dim\operatorname{im}\mathcal C_k=r_k.}
 \tag{PC34}
\]

To display all collision data, at an ambient sum $\lambda$ write
$x=S-\lambda$ and
$\chi_{c,k}(S)=x^{r_\lambda}a_\lambda(x)$ with
$a_\lambda(0)\ne0$.  The kernel in that original primary factor is
the image of the explicit map
\[
 \mathbb C[x]/(x^{m_\lambda-r_\lambda})
  \longrightarrow\mathbb C[x]/(x^{m_\lambda}),\qquad
 [b]\longmapsto[x^{r_\lambda}a_\lambda(x)b].
 \tag{PC35}
\]
It is injective: multiply by the recursively computed inverse of
$a_\lambda$ and cancel the displayed power of $x$.
It is onto the local ideal by its definition.  Thus all
$m_\lambda-r_\lambda$ kernel jets remain, with action
$S=\lambda+x$ and the local unit $a_\lambda$ retained.
If $\Re\lambda\ne k/2$, no all-critical tuple has that sum,
so $r_\lambda=0$ and the entire block belongs to the kernel.
At a sum with real part $k/2$, the same exact formula still
applies: an absent critical tuple kills the whole block, and a
coinciding critical tuple retains precisely the quotient of length
$r_\lambda$ while the additional jets lie in (PC35).
This is the complete relationship between factorwise critical
selection and selection by the real part of the sum.

The original sum action also descends with its exact type.  Multiplication
by $1/2+iy$ is continuous on the one-factor quotient and preserves its
relation space.  On the completed tensor target define the continuous
operator
\[
 M^{[k]}=\sum_{a=1}^k
       I^{\widehat\otimes(a-1)}\widehat\otimes M
                         \widehat\otimes I^{\widehat\otimes(k-a)}.
 \tag{PC36}
\]
On pure tensors this is the sum of the original $k$ multipliers,
and continuity extends it to the completed target.  Since each
$\eta_h$ commutes with the original one-factor action, (PC12)
gives $\mathcal C_kA_k=M^{[k]}\mathcal C_k$ on the cyclic image.
At source level, the finite multinomial expansion of
$P(D_1+\cdots+D_k)F_h^{\otimes k}$ and (PC6) give
$\mathsf q^{\widehat\otimes k}\mathcal V_kP
=\sigma_h^{\widehat\otimes k}\eta_k[P]$.
Thus both the completed comparison and its sum action are induced
by the original weighted source, with no replacement of $U_k$.

\subsection{The same kernel curvature in every original cyclic tensor packet}

Form the original period matrix $\Pi_k(t)$ of (SC1)--(SC6) with
the literal polynomial $\chi_{h,k}$, its full increasing-power basis,
the same specified $u\ne0$, and its own degree $q_k$ in the ray
formula.  Write $e_{0,k}=[1]$,
$\ell_k[P]=[S^{q_k-1}]\operatorname{rem}_{\chi_{h,k}}P$ and
$A_{k,t}=A_k+t e_{0,k}\ell_k$.  The preceding convergence,
remainder and determinant argument applies to this polynomial with
all its coefficients; hence $\Pi_k'=-\Pi_kA_{k,t}/u$ and it is
invertible at every $t$.  Define
\[
 \mathcal C_{k,t}=\mathcal C_k\Pi_k^{-1},\quad
 A_{k,t}^{\rm per}=\Pi_kA_{k,t}\Pi_k^{-1},\quad
 \mathcal K_{k,t}=\Pi_k F_k.
 \tag{PC37}
\]
Equations (PC34) and (PC36) prove
\[
 \ker\mathcal C_{k,t}=\mathcal K_{k,t},\qquad
 \mathcal C_{k,t}A_{k,t}^{\rm per}-M^{[k]}\mathcal C_{k,t}
       =t\mathcal C_k(e_{0,k})\ell_k\Pi_k^{-1}.
 \tag{PC38}
\]
If $0<r_k<q_k$, then $\mathcal C_k(e_{0,k})\ne0$ by
(PC31)--(PC33).  Put $p_k=q_k-r_k$.  The kernel ideal has basis
$[\chi_{c,k}],\ldots,[S^{p_k-1}\chi_{c,k}]$, with the last
member monic of degree $q_k-1$.  Consequently $\ell_k|F_k$ is
onto, its kernel has dimension $p_k-1$, and $e_{0,k}\notin F_k$.
For every $t\ne0$, the restriction of (PC38) to $\mathcal K_{k,t}$
has rank one, image $\mathbb C\mathcal C_k(e_{0,k})$, and full
kernel $\Pi_k(F_k\cap\ker\ell_k)$.  At zero it vanishes.

With any specified inclusion $I_k:F_k\hookrightarrow E_k$, put
$Y_k=\Pi_kI_k$, $H_k=Y_k^*Y_k$,
$N_k=I-Y_kH_k^{-1}Y_k^*$, $z_k=N_k\Pi_ke_{0,k}$ and
$L_k=\ell_kI_k$.  The calculations (PC23)--(PC24) now give
\[
 \boxed{\partial_{\bar t}\partial_t\log\det H_k
   =\frac{|t|^2}{|u|^2}\|z_k\|^2L_kH_k^{-1}L_k^*.}
 \tag{PC39}
\]
Its coefficient is strictly positive for $0<r_k<q_k$, and the
normal acceleration at zero is $-z_k(0)L_k/u$ with rank one.
This calculation includes the collision-kernel jets of (PC35).
The complete neighboring-minor form is obtained with
\[
 \begin{gathered}
 J_g=[\chi_{c,k},S\chi_{c,k},\ldots,S^{p_k-1}\chi_{c,k}],
 \quad J_-=[\chi_{c,k},\ldots,S^{p_k-2}\chi_{c,k}],\\
 J_+=[J_g,e_{0,k}],\quad
 h_J=\det((\Pi_kJ)^*\Pi_kJ).
 \end{gathered}
 \tag{PC40}
\]
For $p_k=1$, $J_-$ is empty and $h_-=1$.  The Schur complement
of the last column gives $h_+=h_g\|z_k\|^2$; the adjugate
formula for the last inverse-Gram entry gives
$L_kH_k^{-1}L_k^*=h_-/h_g$ after the exact constant frame
congruence of (SC32).  Therefore
\[
 \partial_{\bar t}\partial_t\log h_g
       =\frac{|t|^2}{|u|^2}\frac{h_+h_-}{h_g^2}.
 \tag{PC41}
\]
For the original tensor-source Gram $G_{N,k}$, each of these same
inclusions has factors
$v_{N,J}=\det(J^*G_{N,k}J)$ and
$b_{N,J}=\det((J^*G_{N,k}J)^{-1}J^*\Pi_k^*\Pi_kJ)$.
The literal product identity $h_J=v_{N,J}b_{N,J}$ retains every
theta mass, polynomial map and period factor in (PC41).
The moving metric $\Pi_k^{-*}G_{N,k}\Pi_k^{-1}$ has constant
kernel Gram $I_k^*G_{N,k}I_k$, with the full derivative cancellation
(PC27), replacing $A_t$ by $A_{k,t}$.  These are the two exact
metric observations on the same critical-comparison kernel.

If $Z_c=\varnothing$, then $r_k=0$, $F_k=E_k$ and
$\mathcal C_k=0$; the full kernel bundle has zero curvature.
If all packet zeros are critical, $\chi_{c,k}=\chi_{h,k}$,
$F_k=0$, and the empty-kernel curvature is zero.
If both sectors occur, an all-critical tuple proves $r_k>0$,
while a repeated off-critical root has sum $k\rho$ of real part
different from $k/2$, proving $r_k<q_k$.  Thus the strict case
in (PC39) has an exact sector criterion at every original $k$.

Finally each coefficient arrow in this section has its original
Split-Zero lift: $v^\bullet\mapsto(Tv)^\bullet$ and
$\tau\mapsto\tau$, with the retained support label.
The complete kernel (PC34), its smaller action-defect kernel in
(PC38), and its cohomological source image remain separately
represented.  A vector killed by one of these comparisons lands at
that receiving fibre's supported zero.  The amplitude projection
maps both that zero and external absence to ordinary zero, and the
displayed split maps commute with that projection.  This preserves
the original support and the full arithmetic unit through the
critical quotient, the completed tensor target, and the period
curvature calculation.
